%% LyX 2.4.4 created this file.  For more info, see https://www.lyx.org/.
%% Do not edit unless you really know what you are doing.
\documentclass[english]{article}
\usepackage[T1]{fontenc}
\usepackage[latin9]{inputenc}
\usepackage{enumitem}
\usepackage{amsmath}
\usepackage{amsthm}
\usepackage{amssymb}
\usepackage{geometry}
\geometry{verbose,tmargin=1in,bmargin=1in,lmargin=1in,rmargin=1in}

\makeatletter
%%%%%%%%%%%%%%%%%%%%%%%%%%%%%% Textclass specific LaTeX commands.
\numberwithin{equation}{section}
\numberwithin{figure}{section}
\newlength{\lyxlabelwidth}      % auxiliary length 

%%%%%%%%%%%%%%%%%%%%%%%%%%%%%% User specified LaTeX commands.
\usepackage{fancyhdr}
\usepackage{lastpage}
\pagestyle{fancy}
\usepackage{enumitem}
\usepackage{ifthen}
\everymath=\expandafter{\the\everymath\displaystyle}
%\DeclareMathSizes{10}{10}{10}{10}
\usepackage{multicol}

\usepackage{tikz}
\usepackage{tikz-3dplot}
\usetikzlibrary{calc,arrows}

\usetikzlibrary{patterns}
\usetikzlibrary{plotmarks}
\usepackage{pgfplots}
\pgfplotsset{compat=newest}

\usepackage{calc}
\usepackage[nomessages]{fp}

\usepackage{datetime}
% Added by lyx2lyx
\setlength{\parskip}{\smallskipamount}
\setlength{\parindent}{0pt}

\makeatother

\usepackage{babel}
\begin{document}
\renewcommand{\author}{Dr.\,James Ashe}

\newcommand{\classNum}{336}

\newcommand{\classSec}{A}

\newcommand{\examType}{HW 3.3}

\renewcommand{\date}{\formatdate{21}{10}{2013}}

%%First page's header%%%%%%%%%%%%%%%%%%%%%%
\fancypagestyle{firststyle} %%header settings for first pagr
{
	\fancyhead[L]{MTH \classNum{}(\classSec{}) \author{}\\
	\textbf{\examType{}} ~\date{}}
	\fancyhead[C]{}
	\fancyhead[R]{\textbf{Solutions}}
	\setlength{\headsep}{14pt}
}
%%%%%%%%%%%%%%%%%%%%%%%%%%%%%%%%%%%%%%%%%%

%%Next page's header%%%%%%%%%%%%%%%%%%%%%%%
\setlist{nolistsep}
\lhead{\classNum{}(\classSec{})} 
\chead{\textbf{Solutions}} 
\cfoot{\thepage{}~of~\pageref{LastPage}} 
\setlength{\headsep}{5pt} 
%%%%%%%%%%%%%%%%%%%%%%%%%%%%%%%%%%%%%%%%%%%

%%Various settings; see the preamble for other settings%%%%%
%\setenumerate[0]{leftmargin=12pt,itemindent=0pt,labelwidth=0pt}
\setlist{topsep=.5pt}
\renewcommand{\labelenumi}{\textbf{\arabic{enumi}.}}
\renewcommand{\labelenumii}{\textbf{\alph{enumii}.}}
\thispagestyle{firststyle}
%%%%%%%%%%%%%%%%%%%%%%%%%%%%%%%%%%%%%%%%%%

\newcommand{\xmax}{0}
\newcommand{\xmin}{-\xmax}
\newcommand{\xstep}{0}
\newcommand{\ymax}{0}
\newcommand{\ymin}{-\ymax}
\newcommand{\ystep}{0} 
\newcommand{\dspace}{\vspace{1cm}}
\newcommand{\xa}{0}
\newcommand{\xb}{0}
\newcommand{\ya}{1}
\newcommand{\yb}{1}
\begin{enumerate}[resume]
\item \vspace{0cm}

\begin{enumerate}[resume]
\item Correct.
\[
\begin{bmatrix}\begin{array}{rrr}
1 & 2 & 2\\
2 & 4 & 4\\
3 & 5 & 6
\end{array}\end{bmatrix}\rightsquigarrow\begin{bmatrix}\begin{array}{rrr}
1 & 0 & 2\\
0 & 1 & 0\\
0 & 0 & 0
\end{array}\end{bmatrix}
\]
so $2(1,2,3)=(2,4,6)$, meaning that $\mathbf{v}_{2}$ can be written
as a linear combination of $\mathbf{v}_{1}.$ Thus $2\mathbf{v}_{1}-\mathbf{v}_{3}=\mathbf{0}$
is a nontrivial linear combination that equals $\mathbf{0}$.
\item Incorrect. $\mathbf{v}_{2}$ cannot be written as a combination of
$\mathbf{v}_{1}$ and $\mathbf{v}_{3}$. To see this solve the augmented
matrix, 
\[
\begin{bmatrix}\begin{array}{rr|r}
1 & 2 & 2\\
2 & 4 & 4\\
3 & 6 & 5
\end{array}\end{bmatrix}\rightsquigarrow\begin{bmatrix}\begin{array}{rr|r}
1 & 2 & 0\\
0 & 0 & 1\\
0 & 0 & 0
\end{array}\end{bmatrix}
\]
 which is inconsistent. 
\end{enumerate}
\item \vspace{0cm}

\begin{enumerate}
\item Yes. Suppose that $c_{1}\begin{bmatrix}\begin{array}{r}
1\\
4
\end{array}\end{bmatrix}+c_{2}\begin{bmatrix}\begin{array}{r}
2\\
9
\end{array}\end{bmatrix}=\mathbf{0}$
\[
\begin{bmatrix}\begin{array}{rr}
1 & 2\\
4 & 9
\end{array}\end{bmatrix}\rightsquigarrow\begin{array}{c}
\\\mbox{R2-4R1}
\end{array}\begin{bmatrix}\begin{array}{rr}
1 & 2\\
0 & 1
\end{array}\end{bmatrix}\rightsquigarrow\begin{bmatrix}\begin{array}{rr}
1 & 0\\
0 & 1
\end{array}\end{bmatrix}
\]
 implies that $c_{1}=c_{2}=0$. Thus the vectors are linearly independent.
\end{enumerate}
\begin{enumerate}[resume]
\item [\textbf{d.}]As above we need to put the associated column matrix
into reduced echelon form.
\[
\begin{bmatrix}\begin{array}{rrr}
1 & 2 & 0\\
1 & 3 & 1\\
1 & 3 & 2
\end{array}\end{bmatrix}\rightsquigarrow\begin{bmatrix}\begin{array}{rrr}
1 & 2 & 0\\
1 & 3 & 1\\
0 & 0 & 1
\end{array}\end{bmatrix}\rightsquigarrow\begin{bmatrix}\begin{array}{rrr}
1 & 2 & 0\\
1 & 3 & 0\\
0 & 0 & 1
\end{array}\end{bmatrix}\rightsquigarrow\begin{bmatrix}\begin{array}{rrr}
1 & 2 & 0\\
0 & 1 & 0\\
0 & 0 & 1
\end{array}\end{bmatrix}\rightsquigarrow\begin{bmatrix}\begin{array}{rrr}
1 & 0 & 0\\
0 & 1 & 0\\
0 & 0 & 1
\end{array}\end{bmatrix}
\]
so the vectors are linearly independent. Note that you only need to
reduce the matrix to 
\[
\begin{bmatrix}\begin{array}{rrr}
a_{11} & a_{12} & a_{13}\\
0 & a_{22} & a_{23}\\
0 & 0 & a_{33}
\end{array}\end{bmatrix}
\]
where it is clear the matrix has three pivots and thus a rank of $3$.
\end{enumerate}
\item The span of the vectors need to be equal to the given set, and the
set needs to be linearly independent. 

\begin{enumerate}[resume]
\item [\textbf{b.}]No.\\
\[
\begin{bmatrix}\begin{array}{rrrr}
1 & 1 & 2 & 2\\
0 & 2 & 2 & 2\\
1 & 4 & 5 & -1
\end{array}\end{bmatrix}\rightsquigarrow\begin{bmatrix}\begin{array}{rrrr}
1 & 0 & 1 & 0\\
0 & 1 & 1 & 0\\
0 & 0 & 0 & 1
\end{array}\end{bmatrix}
\]
So there is a non-trivial $c_{1},c_{2},c_{3},c_{4}$ such that $c_{1}(1,0,1)+c_{2}(1,2,4)+c_{3}(2,2,5)+c_{4}(2,2,-1)=\mathbf{0}$.
\\
To find it, $c_{1}=-c_{3}$, $c_{2}=-c_{3}$, and $c_{4}=0$. For
$c_{3}=-1$ then $1(1,0,1)+1(1,2,4)-(2,2,5)+0(2,2,-1)=\mathbf{0}$.\\
Alternatively, by proposition 3.4 if the column vectors from $\begin{bmatrix}\begin{array}{rrrr}
1 & 1 & 2 & 2\\
0 & 2 & 2 & 2\\
1 & 4 & 5 & -1
\end{array}\end{bmatrix}$ are a basis of $\mathbb{R}^{3}$ then it is a $3\times3$ nonsingular
matirx but it is a $3\times4$.

\begin{enumerate}[resume]
\item []Any four vectors in $\mathbb{R}^{3}$ form a linearly dependent
set. In general, any $n$ vectors in $\mathbb{R}^{n-1}$ (for $n>1$)
form a linearly dependent set. This is a result of proposition 4.1
or 3.4, but you can see this by evaluating the possible reduced echelon
forms of the augmented column matrix, 
\[
A\mathbf{x}=\mathbf{b}=\begin{bmatrix}\begin{array}{rrr|r}
a_{11} & \cdots & a_{1n} & a_{1,n+1}\\
\vdots &  & \vdots & \vdots\\
a_{n1} & \cdots & a_{nn} & a_{n,n+1}
\end{array}\end{bmatrix}\mbox{.}
\]
An inconsistent matrix means the vectors composing $A$ are linearly
dependent, but a solvable matrix means that $\mathbf{b}$ can be written
as a linear combination of the vectors composing $A$. Either way
the set of vectors is linearly dependent.
\end{enumerate}
\item [\textbf{c.}]No. 
\[
\begin{bmatrix}\begin{array}{rrr}
1 & 0 & 1\\
0 & 1 & 1\\
2 & 1 & 4\\
3 & 1 & 4
\end{array}\end{bmatrix}\rightsquigarrow\begin{bmatrix}\begin{array}{rrr}
1 & 0 & 0\\
0 & 1 & 0\\
0 & 0 & 1\\
0 & 0 & 0
\end{array}\end{bmatrix}
\]
Although these vectors are linearly independent. The span of $3$
vectors has a dimension of at most $3$. For example, $(0,0,0,1)$
is not in the span of these vectors since 
\[
\begin{bmatrix}\begin{array}{rrr|}
1 & 0 & 1\\
0 & 1 & 1\\
2 & 1 & 4\\
3 & 1 & 4
\end{array}\begin{array}{r}
0\\
0\\
0\\
1
\end{array}\end{bmatrix}\rightsquigarrow\begin{bmatrix}\begin{array}{rrr}
1 & 0 & 0\\
0 & 1 & 0\\
0 & 0 & 1\\
0 & 0 & 0
\end{array}\begin{array}{|r}
0\\
0\\
0\\
1
\end{array}\end{bmatrix}
\]
is inconsistent. 
\end{enumerate}
\item [\textbf{8.}]Let $c_{1}(\mathbf{v}-\mathbf{w})+c_{2}(2\mathbf{v}+\mathbf{w})=\mathbf{0}$.
(We need to show that $c_{1}=c_{2}=0$ which implies that the vectors
are linearly independent).
\begin{eqnarray*}
c_{1}(\mathbf{v}-\mathbf{w})+c_{2}(2\mathbf{v}+\mathbf{w}) & = & \mathbf{0}\\
\Rightarrow\mathbf{v}(c_{1}+2c_{2})+\mathbf{w}(c_{2}-c_{1}) & = & \mathbf{0}
\end{eqnarray*}
Since $\{\mathbf{v},\mathbf{w}\}$ is linearly independent, $c_{1}+2c_{2}=0$
and $c_{2}-c_{1}=0$. This gives the homogeneous system of equations,
\[
\begin{bmatrix}\begin{array}{rr}
1 & 2\\
-1 & 1
\end{array}\end{bmatrix}\rightsquigarrow\begin{bmatrix}\begin{array}{rr}
1 & 0\\
0 & 1
\end{array}\end{bmatrix}
\]
 which implies that $c_{1}=c_{2}=0$ as needed,
\end{enumerate}

\end{document}
